\section{Proposed Method: CAC-CycleGAN-WGP}
\label{sec:proposed_method}

This section details the proposed CAC-CycleGAN-WGP framework for intelligent fault diagnosis under imbalanced data conditions. We describe the overall framework, the specific architectural design of the generative networks, the formulation of the objective functions, and the subsequent quality evaluation and diagnosis stages.

\subsection{Overall Framework and Operational Flow}
The proposed CAC-CycleGAN-WGP method aims to synthesize high-quality faulty samples by learning the bidirectional mapping between the normal signal domain and various fault signal domains. Unlike traditional GANs that generate samples from random noise, our approach leverages the inherent structural information of normal vibration signals to "transform" them into realistic faulty signals. The overall operational flow, as illustrated in Fig. \ref{fig:flowchart}, consists of five primary stages:

\begin{enumerate}
    \item \textbf{Data Preprocessing:} Vibration signals collected from sensors are segmented using a sliding window of 1024 points. Fast Fourier Transform (FFT) or time-domain normalization is applied to convert raw temporal sequences into feature-rich representations suitable for deep learning architectures.
    \item \textbf{CycleGAN Training:} The preprocessed normal and faulty samples are used to train the dual-generator framework. The CAC-CycleGAN-WGP model optimizes the mapping functions through a combination of adversarial, cycle-consistency, and auxiliary classification losses, stabilized by Wasserstein distance and gradient penalty.
    \item \textbf{Sample Generation:} Once the generators reach convergence, the generator $G: X \to Y$ is employed to transform healthy samples into synthetic faulty samples for all minority classes.
    \item \textbf{Quality Evaluation (SAE):} A Stacked Autoencoder (SAE) is utilized to evaluate the physical fidelity and diversity of the generated signals. Quantitative metrics such as Pearson Correlation Coefficient (PCC) and Cosine Similarity (CS) are used to filter out low-quality synthetic data.
    \item \textbf{Fault Diagnosis (SVM):} The high-quality synthetic samples are merged with the original imbalanced training set to form a balanced/augmented dataset. A Support Vector Machine (SVM) classifier is trained on this augmented set and validated on a separate, real-world test set.
\end{enumerate}

By bridging the gap between healthy and faulty distributions through cycle-consistency, the model ensures that the generated signals retain the fundamental mechanical characteristics of the bearing while exhibiting the discriminative features of specific faults.

\subsection{Architecture Design}
The effectiveness of CAC-CycleGAN-WGP relies on a sophisticated deep CNN-based architecture, comprising two generators, two discriminators, and a Conditional Auxiliary Classifier (CAC). The design is specifically tailored to handle the high-dimensional and periodic nature of raw vibration signals, ensuring that both global temporal structures and local transient features (impulse components) are preserved during the domain translation process.

\subsubsection{Generator Architecture}
The generators ($G_{X \to Y}$ and $F_{Y \to X}$) adopt a symmetrical encoder-decoder structure, which is essential for capturing and reconstructing the complex features of vibration data. The encoder consists of three 1D convolutional layers with filter counts of 64, 128, and 256. For an input signal of length 1024, the first layer uses a large kernel size of 7 with stride 1 to capture long-range dependencies and filter out initial high-frequency white noise. The subsequent layers use a kernel size of 3 and stride 2 for downsampling, effectively compressing the 1024-point vector into a 256-dimensional latent space. 

The core of the generator incorporates six residual blocks (ResBlocks) to prevent gradient degradation during deep feature transformation. Each ResBlock contains two 1D convolutional layers with 256 filters and a kernel size of 3. Residual learning allows the model to learn the "residual" difference between the healthy and faulty states, which is mathematically more efficient than learning the entire target distribution from scratch. This bottleneck design forces the network to focus on the periodic impact features that differentiate a faulty bearing from a healthy one.

The decoder utilizes two 1D transposed convolutions (deconvolution) with stride 2 to upsample the latent representation back to its original length of 1024 points. Instance normalization is applied after each convolutional layer instead of Batch Normalization to avoid artifacts related to the batch size, ensuring that the statistics of individual vibration signals are preserved. ReLU activation functions are employed throughout the hidden layers to maintain sparsity, while a Tanh activation at the final output layer scales the signals within $[-1, 1]$. This specific architecture ensures that the generator can effectively capture the multi-scale periodic patterns characteristic of mechanical vibrations, ranging from low-frequency modulation to high-frequency resonance.

\subsubsection{Discriminator and CAC Architecture}
The discriminators ($D_X$ and $D_Y$) are designed as PatchGAN-style convolutional networks. Unlike traditional discriminators that output a single scalar for the entire input, the PatchGAN structure classifies local segments (patches) of the signal as real or fake. This approach encourages the generator to produce high-frequency details and local consistency across the entire signal length, which is crucial for simulating the localized impacts produced by bearing inner or outer race defects. The discriminators consist of four 1D convolutional layers with 64, 128, 256, and 512 filters, all using a kernel size of 4 and stride 2. This progressively increases the receptive field while reducing temporal resolution, allowing the critic to evaluate the signal at multiple scales.

To incorporate class-specific information and handle the multiclass imbalance problem, the model integrates a Conditional Auxiliary Classifier (CAC). The CAC shares the feature extraction layers of the discriminator but terminates in a separate Softmax branch with $N$ nodes (where $N$ is the number of fault classes). This multi-task configuration serves two purposes: first, it provides a "regularization" effect that helps the discriminator learn more robust features; second, it allows the generator to receive explicit feedback on whether its synthetic output matches the intended fault category. By optimizing the classification accuracy of both real and synthetic samples, the generator is forced to produce signals that are class-distinctive and reside on the target fault manifold. The integration of CAC transforms the standard unsupervised CycleGAN into a semi-supervised framework capable of generating specific types of faults (e.g., ball fault, inner race fault) with high precision, solve the limitation of standard CycleGAN which only performs domain-to-domain translation without class awareness.

\subsection{Objective Functions and Formulation}
The training of CAC-CycleGAN-WGP is governed by a multi-objective loss function. To ensure stability and handle multiclass fault categories, we formulate the optimization problem as follows.

\subsubsection{Adversarial Loss with Wasserstein Distance}
To improve training stability and provide a more informative gradient even when distributions are disjoint, we replace the standard JS-divergence with the Wasserstein distance. The adversarial loss for the mapping $G: X \to Y$ and its critic $D_Y$ is defined as the Earth Mover's distance:
\begin{equation}
\mathcal{L}_{\text{adv}}(G, D_Y, X, Y) = \mathbb{E}_{\boldsymbol{y} \sim p_{data}(\boldsymbol{y})}[D_Y(\boldsymbol{y})] - \mathbb{E}_{\boldsymbol{x} \sim p_{data}(\boldsymbol{x})}[D_Y(G(\boldsymbol{x}))]
\end{equation}
The Wasserstein distance provides a smooth gradient relative to the generator's parameters, which is critical for learning the sparse distributions of fault signals where traditional GANs often fail to converge.

\subsubsection{Cycle-Consistency Loss}
To constrain the mapping and ensure that the generator learns a deterministic transformation rather than random mapping (i.e., avoiding mode collapse), the cycle-consistency loss is enforced. This loss ensures that $F(G(\boldsymbol{x})) \approx \boldsymbol{x}$ and $G(F(\boldsymbol{y})) \approx \boldsymbol{y}$:
\begin{equation}
\mathcal{L}_{\text{cyc}}(G, F) = \mathbb{E}_{\boldsymbol{x} \sim p_{data}(\boldsymbol{x})} [\|F(G(\boldsymbol{x})) - \boldsymbol{x}\|_1] + \mathbb{E}_{\boldsymbol{y} \sim p_{data}(\boldsymbol{y})} [\|G(F(\boldsymbol{y})) - \boldsymbol{y}\|_1]
\end{equation}
The usage of the $L1$ norm encourages the preservation of the temporal structure and background noise characteristics of the source signal, while the adversarial loss injects the fault-specific signatures.

\subsubsection{Conditional Auxiliary Classification Loss}
To ensure the generated samples belong to the target fault category, the CAC loss is added. It is defined as the cross-entropy loss for class prediction. It consists of two parts: the classification loss of real samples ($\mathcal{L}_{C}^{r}$) and generated samples ($\mathcal{L}_{C}^{g}$):
\begin{equation}
\mathcal{L}_{C} = \mathbb{E}[\log P(C=c | \boldsymbol{x}_{real})] + \mathbb{E}[\log P(C=c | G(\boldsymbol{x}))]
\end{equation}
where $P(C| \cdot)$ is the probability distribution over class labels predicted by the CAC. This term forces the generator to synthesize samples that exhibit the discriminative features of the target class $c$, effectively solving the multiclass generation problem within a single framework.

\subsubsection{Gradient Penalty}
To strictly enforce the 1-Lipschitz constraint required by the WGAN framework and avoid the capacity issues associated with weight clipping, a gradient penalty (GP) term is applied to the discriminator:
\begin{equation}
\mathcal{L}_{\text{GP}} = \lambda \mathbb{E}_{\hat{\boldsymbol{x}} \sim p_{\hat{\boldsymbol{x}}}} [(\|\nabla_{\hat{\boldsymbol{x}}} D_Y(\hat{\boldsymbol{x}})\|_2 - 1)^2]
\end{equation}
where $\hat{\boldsymbol{x}}$ is an interpolation between real and fake samples, and $\lambda$ is usually set to 10. This term ensures that the discriminator's gradients do not vanish or explode, providing stable feedback to the generator throughout the training process.

\subsubsection{Total Objective Function}
The final loss function used to optimize the CAC-CycleGAN-WGP model is a weighted sum of the components:
\begin{equation}
\mathcal{L}_{\text{total}} = \mathcal{L}_{\text{adv}}(G, D_Y) + \mathcal{L}_{\text{adv}}(F, D_X) + \alpha \mathcal{L}_{\text{cyc}} + \beta \mathcal{L}_{C} + \gamma \mathcal{L}_{\text{GP}}
\end{equation}
where $\alpha, \beta, \gamma$ are hyperparameters. In our experiments, $\alpha=10$ is used to emphasize structural reconstruction, while $\beta=1.0$ and $\gamma=10$ are set for classification accuracy and training stability. The training algorithm follows an alternating update schedule where the discriminators are optimized for 5 steps for every 1 generator update, ensuring that the critics are close to the optimal Wasserstein distance.

\subsection{SAE-based Quality Evaluation}
To verify that the generated signals are physically meaningful and suitable for diagnosis, we propose an evaluation strategy based on Stacked Autoencoders (SAE). In industrial applications, generating samples is insufficient; they must adhere to the physical laws of machinery dynamics. An SAE is first trained on the original real dataset to learn a compressed, latent feature representation of various fault types. The SAE architecture comprises an input layer (1024), three hidden layers (512, 256, 128) using Sigmoid activation, and a symmetrical decoder. By training to minimize the reconstruction error of real signals, the SAE learns an "expectation" of what a legitimate vibration signal should look like in the feature space.

The evaluation process involves feeding the synthetic samples into the pre-trained SAE. We monitor the reconstruction error and the latent feature distribution. Specifically, we calculate the Pearson Correlation Coefficient (PCC) and Cosine Similarity (CS) between the generated faulty signals ($G(\boldsymbol{x})$) and real faulty signals ($\boldsymbol{y}$) of the same class. 
\begin{equation}
PCC = \frac{\sum (G(\boldsymbol{x})_i - \bar{G(\boldsymbol{x})})(\boldsymbol{y}_i - \bar{\boldsymbol{y}})}{\sqrt{\sum (G(\boldsymbol{x})_i - \bar{G(\boldsymbol{x})})^2 \sum (\boldsymbol{y}_i - \bar{\boldsymbol{y}})^2}}
\end{equation}
PCC measures the linear correlation between the temporal waveforms, ensuring that the phase and amplitude variations are consistent with historical data. Simultaneously, the CS evaluates the directional similarity of the feature vectors in the high-dimensional space.

To implement a rigorous filtering mechanism, a PCC threshold $\tau_{PCC} = 0.85$ and a CS threshold $\tau_{CS} = 0.80$ are empirically set based on the distribution of real samples. Samples failing to meet these thresholds are considered "noisy" or "mode-collapsed" and are discarded from the training pool. Furthermore, we employ a T-distributed Stochastic Neighbor Embedding (t-SNE) visualization to qualitatively check if the synthetic samples are overlapping with the real sample clusters. This dual-verification (quantitative thresholds + qualitative manifold check) ensures that only high-fidelity data is used for training the final classifier, bridging the gap between generative modeling and industrial reliability. It prevents the "garbage in, garbage out" phenomenon that often plagues data augmentation methods in sensitive fields like structural health monitoring.

\subsection{SVM-based Final Diagnosis}
The final stage involves the deployment of a Support Vector Machine (SVM) for fault classification. While deep learning models like CNNs are powerful, SVM is chosen here as the baseline classifier for its robustness in handling high-dimensional data and its effective generalization in small-sample scenarios once the classes are made separable by our generative augmentation. SVM aims to find the optimal hyperplane that maximizes the margin between different fault categories.

After balancing the training set using synthetic samples from CAC-CycleGAN-WGP, we define the Balance Ratio (BR) as the ratio of minority class samples to the majority class samples:
\begin{equation}
BR = \frac{N_{synthetic} + N_{minority}}{N_{majority}}
\end{equation}
By adjusting $BR$ from its initial imbalanced state (e.g., $BR=0.1$) to a balanced state ($BR=1.0$), we provide the classifier with a comprehensive view of the fault manifold. The SVM model utilizes a Radial Basis Function (RBF) kernel to handle nonlinearities in the vibration feature space. The RBF kernel map is given by $K(x, x') = \exp(-\gamma ||x - x'||^2)$, where the kernel parameter $\gamma$ and the penalty parameter $C$ (which controls the trade-off between margin maximization and error minimization) are optimized via 5-fold cross-validation on the augmented training set. 

By training on the augmented dataset, the SVM's decision boundaries are pushed away from the sparse regions where minority samples previously resided. This significantly reduces the empirical risk and bias toward the majority class (usually the "Healthy" condition), improving the sensitivity to rare machinery failures. In the operational phase, the trained SVM receives real-time vibration features and outputs a diagnostic label. The performance is evaluated using the G-mean and F1-score to ensure that the improvement in minority class accuracy does not come at the expense of overall precision. This integration of a state-of-the-art generative model with a classic robust classifier ensures both high accuracy and high reliability in practical industrial environments where data imbalance is the norm rather than the exception.

